% Figure: force on a magnetic dipole in a field gradient. A uniform field
% exerts only a torque; a gradient exerts a translational force that pulls
% a paramagnetic dipole toward high field and pushes a diamagnetic one
% toward low field. The force tracks the field-gradient product B dB/dx.
\begin{figure}[htbp]
\centering
\begin{tikzpicture}[>=Latex]
\begin{axis}[
  width=11.5cm, height=7cm,
  xlabel={position along the gradient $x$ ($\si{\milli\meter}$)},
  ylabel={force per unit moment $F/m$ ($\si{\tesla\per\meter}$)},
  xmin=0, xmax=20,
  ymin=-40, ymax=140,
  xtick={0,4,8,12,16,20},
  legend pos=north east,
  legend cell align=left,
  font=\scriptsize,
  grid=both,
  grid style={enggray!20},
]
% Model field B(x) = B0 exp(-x/L), so dB/dx = -(B0/L) exp(-x/L). The force
% on a moment aligned with the field is m dB/dx (magnitude m|dB/dx|). We
% plot the gradient magnitude |dB/dx| = (B0/L) exp(-x/L) with B0=0.6 T,
% L=8 mm: |dB/dx| at x=0 is 0.6/0.008 = 75 T/m. Curve in T/m.
\addplot[engblue, line width=1.2pt, domain=0:20, samples=120]
  {75*exp(-x/8)};
\addlegendentry{paramagnetic: pulled toward high field}
% diamagnetic sign-reversed, smaller magnitude (weaker susceptibility)
\addplot[engred, line width=1pt, domain=0:20, samples=120]
  {-8*exp(-x/8)};
\addlegendentry{diamagnetic: pushed toward low field}
% mark the high-field edge where the separator gate is placed
\addplot[enggray, only marks, mark=*, mark size=2pt]
  coordinates {(2,58.4)};
\node[enggray, anchor=west, font=\scriptsize] at (axis cs:2.3,66)
  {capture zone};
\end{axis}
\end{tikzpicture}
\caption[Force on a magnetic dipole in a field gradient]{Force per unit
magnetic moment on a small dipole placed in a decaying field
$B(x) = B_0 e^{-x/L}$, with $B_0 = 0.6\,\si{\tesla}$ and decay length
$L = 8\,\si{\milli\meter}$. A uniform field exerts only a torque; the
translational force needs a gradient and equals $m\,\dd B/\dd x$ for a
moment aligned with the field. A paramagnetic particle acquires a moment
along the field and is drawn toward the high-field edge, the mechanism of
magnetic separation and of the gradient coils that localise an MRI signal.
A diamagnetic particle acquires a moment against the field and is pushed
the other way, weakly, because its susceptibility is smaller by orders of
magnitude. The capture zone sits where the gradient is strongest, at the
pole edge, not at the field maximum.}
\label{fig:vol04ch08:dipole-gradient-force}
\end{figure}
